\documentclass[10pt]{article}

% Palatino
\usepackage{mathpazo}

% larger line spacing
\usepackage{setspace}\setstretch{1.15}

% larger margins
\usepackage[margin=1in]{geometry}

% better quotations
\usepackage{csquotes}

% todo notes
\usepackage{todonotes}\setuptodonotes{inline}

% better bullets
\usepackage{enumitem}
\renewcommand{\labelitemi}{$\triangleright$}

% general math
\usepackage{amsthm, amsmath, amssymb, mathtools}
\usepackage{thmtools, thm-restate}

% better links
\usepackage{hyperref}
\usepackage{cleveref}

\title{
  Higher Category Theory via Rzk
}
\author{Directed Study}
\date{Fall 2026}

\begin{document}

\maketitle
%% \tableofcontents

\abstract{
  This is a document for organizing ideas and material for a directed study on higher category theory using the Rzk proof assistant \cite{rzk}.
  The primary goal of this direct study (for me) is to understand to what extend the study of higher category theory at an introductory level can be made simpler or more interesting by using a proof assistant designed for higher category theory.
  Figure \ref{fig:schedule} has a rough schedule for the semester.
  \begin{figure}
    \begin{tabular}{|c|l|}
      \hline
      Week & Topic \\
      \hline
      1 & None \\
      2 & None \\
      3 & Getting set up \cite{rzk} \\
      4 & Introduction to homotopy type theory (\cite{weinberger-2026}, Lecture 1) \\
      5 & Simplicial homotopy and Segal types (\cite{weinberger-2026}, Lecture 2) \\
      6 & Category theory with Segal types (\cite{weinberger-2026}, Lecture 3) \\
      7 & Discrete and Rezk types (\cite{weinberger-2026}, Lecture 4) \\
      8 & Basic higher category theory, (\cite{haugseng-2025}, \S 1) \\
      9 & $\infty$-categories, (\cite{haugseng-2025}, \S 2) \\
      10 & $(\infty, n)$-categories, (\cite{haugseng-2025}, \S 3) \\
      11 & Higher Yoneda's Lemma \\
      12 & TBD \\
      13 & TBD \\
      14 & TBD \\
      15 & TBD \\
      \hline
    \end{tabular}
    \centering
    \caption{An outline of the topics covered during the semester}
    \label{fig:schedule}
  \end{figure}
}

\section{Getting step up}

The goal of this week is to make sure we're all set up to work with Rzk.
I installed Rzk on my machine, there were some issues building via \texttt{cabal} having to do with the \texttt{lsp-types} package, so I gave up and installed the binary from Github.
I played around very breifly with it, the features related to dependent types seem standard, but I haven't dealt with cubes or topes yet.

\subsection*{Exercises}

\enquote{Due} Friday September 25.

\begin{enumerate}
\item Read through Lecture 1 of \textit{Synthetic $\infty$-category theory} \cite{weinberger-2026}.
\item Choose your favorite basic theorem of arithmetic and prove it in Rzk.
\item Formalize Eckmann-Hilton (\cite{hottbook}, Theorem 2.1.6) in Rzk.
\end{enumerate}

\section{Introduction to homotopy type theory}

\begin{itemize}
\item
  This week we're keeping it simple (at least on paper).
  We're looking at synthetic homotopy theory in Rzk.
  This is covered in the \enquote{Day 2} section.
\item
  \enquote{\dots Kan complexes present the \enquote{standard model} of HoTT: a type is a Kan complex, a.k.a.\ $\infty$-groupoid, a.k.a.\ a topological space up to homotopy.}
  According to nLab \cite{nlab-kan-complex}, \enquote{a \textit{Kan complex} is an abstraction of the combinatorial structure found in the nerve of a groupoid.}
  I don't know what this means, and clicking around nLab didn't help.
  Hopefully they'll make it more clear in these notes.
\item
  It took me a moment to realize why we call a proof pointwise equality a homotopy.
  A homotopy from $f : X \to Y$ to $g : X \to Y$ is typically defined as a continuous function $h : X \times [0, 1] \to Y$ such that $h(\cdot, 0) = f$ and $h(\cdot, 1) = g$.
  For any $x \in X$, the function $h(x, \cdot) : [0, 1] \to Y$ is a path from $f(x)$ to $g(x)$.
  In the Curry-Howard-Voevodsky correspondense, paths are equality types.
  The \enquote{Curried} version of $h$ is something that maps $x : X$ to paths (equalities) between $f(x)$ and $g(x)$ (seems obvious in retrospect).
\item
  There is a typo (as of 2026-09-21) in the definition of $\textsf{is-equiv}_{A, B}(f)$.
  It should read:
  \begin{displaymath}
    \left(\sum_{r^{B \to A}}r \circ f \sim i_A \right) \times \left(\sum_{s^{B \to A}} (f \circ s \sim i_B)\right)
  \end{displaymath}
\end{itemize}




\nocite{*}
\bibliographystyle{plain}
\bibliography{refs}

\end{document}
